% LTeX: language=en-EN

\section{Linear Parabolic Problems}

For this section we loosely follow \cite[Chapter 11]{QuarteroniValli1994}.
Parabolic problems like the heat equation can be written as
\begin{equation*}
    \frac{\diff u}{\diff t}(t) + A u(t) = f(t), \qquad \forall t \in (0,T).
\end{equation*}
If $\prodH{Au}{v} = a(u,v)$ for $u \in D_H(A)$ and $v \in V$ (see \Cref{sec:Galerkin linear problems Hilbert space}) we arrive at the weak formulation
\begin{equation}
    \tag{C}
    \label{eq:Galerkin parabolic variational}
    \text{Find }u \in C([0,T]; V) \cap C^1([0,T]; H)\text{ such that }
    \begin{dcases}
        \prodH{\frac{\diff u}{\diff t} (t)}{v} + a(u(t),v) = \prodH{f(t)}{v} , \\
        \qquad   \forall t \in (0,T), v \in V,                        \\
        u(0) = u_0.
    \end{dcases}
\end{equation}

A good presentation of results on existence, uniqueness, and properties for this problem can be found in \cite{Brezis2010}.

\subsection{Energy estimates for \eqref{eq:Galerkin parabolic variational}}

\begin{lemma}[Parabolic a priori estimates]
    \label{lem:Galerkin parabolic a priori estimate}
    Let $a: V \times V \to \mathbb R$ be bilinear, symmetric, continuous and coercive with constant $\alpha^*$, assume that $\|v\|_H \le \|v\|_V$ for all $v \in V$, and let $u$ be a solution to \eqref{eq:Galerkin parabolic variational}. Then we have the estimate
    \begin{equation*}
        \|u(t)\|_H^2 + \int_0^t \|u(s)\|_a^2 \diff s \le \|u_0\|_H^2 +  \frac{1}{\alpha^*} \int_0^t \|f(s)\|_H^2 \diff s, \qquad \forall t \in (0,T).
    \end{equation*}
    Furthermore, we have that
    \begin{equation*}
        \|u(t)\|_H^2 \le e^{-\alpha^* t} \|u_0\|_H^2 + \frac{1}{\alpha^*} \int_0^t e^{-\alpha^*(t-s)} \|f(s)\|_H^2 \diff s.
    \end{equation*}
    If $\frac{\diff u}{\diff t} \in C([0,T]; V)$, then we also have the a priori estimate
    \begin{equation*}
        \int_0^t \left\| \frac{\diff u}{\diff t} (s)\right\|_H^2 \diff s + \|u(t)\|_a^2 \le \|u_0\|_a^2 + \int_0^t \|f(s)\|_H^2 \diff s.
    \end{equation*}
\end{lemma}
\begin{proof}
    Taking $v = u(t)$ and applying the Cauchy--Schwarz and Young inequalities we recover that
    \begin{equation*}
        \frac 1 2 \frac{\diff}{\diff t} \|u(t)\|_H^2 + a(u(t), u(t)) \le \frac 1 {2\alpha^*} \|f(t)\|_H^2 + \frac {\alpha^*} 2 \|u(t)\|_H^2.
    \end{equation*}
    The first estimate follows from applying coercivity (so that $\alpha^* \|u\|_H^2 \le \alpha^* \|u\|_V^2 \le \|u\|_a^2$) and integrating in time directly.
    The second estimate follows from using coercivity and the integrating factor $e^{\alpha^* t}$.
    The third estimate follows from using as test function $v = \frac{\diff u}{\diff t}(t)$, noticing that by symmetry $a(u, \frac{\diff u}{\diff t}) = \frac 1 2 \frac{\diff}{\diff t} \|u\|_a^2$, and arguing similarly.
\end{proof}

\subsection{Space semi-discretisation}
The corresponding discretisation in ``space'' can be written as
\begin{equation}
    \tag{C$_h$}
    \label{eq:Galerkin parabolic variational space-discrete}
    \text{Find }u_h \in C^1([0,T]; V_h) \text{ such that }
    \begin{dcases}
        \prodH{\frac{\diff u_h}{\diff t} (t)}{v_h} + a(u_h(t),v_h)  = \prodH{f(t)}{v_h} \\
        \qquad  \forall t \in (0,T), v_h \in V_h,                              \\
        u_h(0) = u_{0,h}.
    \end{dcases}
\end{equation}

\begin{strongtheorem}[Parabolic version of Céa's lemma]
    \label{thm:Galerkin parabolic Cea}
    This result is a small generalisation of the type of result in \cite[Chapter 11]{QuarteroniValli1994}.
    Under the assumptions of \Cref{lem:Galerkin parabolic a priori estimate}, let $u$ and $u_h$ be the solutions to \eqref{eq:Galerkin parabolic variational} and \eqref{eq:Galerkin parabolic variational space-discrete}, respectively. Then, we have that
    \begin{equation}
        \begin{aligned}
             & \| u(t) - u_h(t) \|_H^2  + \int_0^t \|u(s) - u_h(s)\|_a^2 \diff s                                                                                \\
             & \qquad \le \|u_0 - u_{0,h}\|_H^2 + \inf_{v_h \in L^2(0,T;V_h)} \int_0^t \Big( C_h(s) \| u(s) - v_h(s) \|_{H} + \| u(s) - v_h(s) \|_{a}^2 \Big) \diff s
        \end{aligned}
    \end{equation}
    where
    \begin{equation*}
        C_h(t) = 2 \sup_{s \in (0,t)} \left( \left\| \frac{\diff u}{\diff t} (s) \right\|_H + \left\| \frac{\diff u_h}{\diff t} (s) \right\|_H \right)
    \end{equation*}
    is finite under the stated time-continuity assumptions. A bound independent of $h$ requires additional regularity; the energy estimate for $u_t$ in $L^2(0,T;H)$ alone does not bound this supremum.
\end{strongtheorem}

\begin{proof}
    Let $e_h = u - u_h$ and let $v_h \in C([0,T] , V_h)$. Subtracting \eqref{eq:Galerkin parabolic variational space-discrete} from \eqref{eq:Galerkin parabolic variational} we have the Galerkin orthogonality $\prodH{\frac{\diff e_h}{\diff t}}{w_h} + a(e_h, w_h) = 0$ for all $w_h \in V_h$. Writing $e_h = (u - v_h) + (v_h - u_h)$ and using this with $w_h = v_h - u_h$, we deduce
    \begin{align*}
        \frac 1 2 \frac{\diff}{\diff t} \prodH{e_h}{e_h} + a(e_h, e_h)
         & = \prodH{\frac{\diff e_h}{\diff t}}{u - v_h} + a(e_h, u - v_h)                                                      \\
         & \le \left\| \frac{\diff e_h}{\diff t} \right\|_H \|u - v_h\|_H + \frac 1 2 \|e_h\|_a^2 + \frac 1 2 \|u - v_h\|_a^2,
    \end{align*}
    where we used the Cauchy--Schwarz inequality for the symmetric form $a$.
    Hence, since $2 \| \frac{\diff e_h}{\diff t}(t) \|_H \le C_h(t)$,
    \begin{equation*}
        \frac{\diff}{\diff t} \prodH{e_h}{e_h} + a(e_h, e_h) \le C_h(t) \|u - v_h\|_H + \|u - v_h\|_a^2.
    \end{equation*}
    We can now integrate in time, and then take the infimum in $v_h \in C([0,T]; V_h)$. The infimum coincides with the infimum over $L^2(0,T; V_h)$ by density.
\end{proof}

\subsubsection{On the choice of $u_{h}(0)$}

There are many possible choices of $u_h(0)$. A convenient one for the analysis is the so-called elliptic (or Ritz) projection of $u_0$, obtained by solving
\begin{equation*}
    u_h(0) \in V_h \text{ such that} \quad a(u_h(0), v_h) = a(u_0, v_h) \quad \forall v_h \in V_h.
\end{equation*}
If $u_0 \in D_H(A)$, then we have that
\begin{equation*}
    \prodH {\frac{\diff u_h}{\diff t}(0)}{v_h} = -a(u_0, v_h) + \prodH{f(0)}{v_h} = \prodH{-Au_0 + f(0)}{v_h} \le \left( \| Au_0 \|_H  + \|f(0) \|_H \right) \|v_h\|_H.
\end{equation*}
Taking $v_h = \frac{\diff u_h}{\diff t}(0)$ we recover that
\begin{equation}
    \label{eq:Galerkin parabolic estimate duhdt(0)}
    \left\| \frac{\diff u_h}{\diff t}(0) \right\|_H \le \| Au_0 \|_H + \|f(0)\|_H.
\end{equation}


If, in addition, $f\in H^1(0,T;H)$ and the differentiated equations admit the energy estimate of \Cref{lem:Galerkin parabolic a priori estimate}, apply that estimate to $u_t$ and $u_{h,t}$. With the Ritz initial value, $u_{h,t}(0)$ is the $H$-orthogonal projection of $f(0)-Au_0$. Thus
\begin{equation*}
    C_h(t) \le 4\left(\|f(0)-Au_0\|_H^2+
    \frac{1}{\alpha^*}\int_0^t\|f'(s)\|_H^2\diff s\right)^{1/2}.
\end{equation*}
This gives an $h$-independent bound. The differentiated energy argument is justified directly for smooth solutions and, by approximation, for the corresponding energy solutions.

\subsubsection{System of linear ODEs for the coefficients}
\label{sec:Galerkin parabolic linear ODEs}
We now look for solutions of the form
\begin{equation*}
    u_h(t) = \sum_{j=1}^{N} U_j(t) \varphi_j
\end{equation*}
with $U_j \in C^1([0,T];\mathbb R)$.
This leads to the linear system of ODEs for the basis coefficients
\begin{equation}
    \label{eq:Galerkin parabolic linear algebra}
    \mat M \frac{\diff \vec u}{\diff t} (t) + \mat K \vec u (t) = \vec b (t), \qquad
    \text{where } \mat M = \Big( \prodH{\varphi_j}{\varphi_i} \Big)_{ij},
    \vec u (t) = \left(U_i(t)\right)_i,
    \vec b (t) = \left(\prodH{f(t)}{\varphi_i}\right)_i.
\end{equation}
Here $\mat K$ is the stiffness matrix introduced in \Cref{sec:Galerkin elliptic linear algebra}.
Matrix $\mat M$ is known as the \emph{mass matrix}.

\subsection{Fully discrete problem}
Then we must pick a suitable time discretisation, in order to be able to numerically solve our problem. For example, we can use the $\theta$-method with $\theta \in [0,1]$ ($\theta = 0$ is the explicit Euler method, $\theta = 1$ the implicit Euler method and $\theta = \frac 1 2$ the Crank--Nicolson method): given $u_h^n$ (an approximation of $u_h(n \tau)$) we take $u_h^{n+1}$ (an approximation of $u_h ((n+1)\tau)$) as the solution to
\begin{equation*}
    \prodH{\frac{u_h^{n+1} - u_h^n}{\tau}}{v_h} + a\Big(\theta u_h^{n+1} + (1-\theta) u_h^n ,v_h\Big)
    = \BigprodH{\theta f((n+1)\tau) + (1-\theta) f(n\tau)}{v_h}, \qquad \forall v_h \in V_h
\end{equation*}
that can be written as the elliptic problem
\begin{equation*}
    \prodH{u_h^{n+1}} {v_h} + \tau \theta a(u_h^{n+1},v_h)
    = \BigprodH{u_h^n + \tau \theta f((n+1)\tau) + \tau (1-\theta) f(n\tau)}{v_h} - \tau (1-\theta) a(u_h^n , v_h) , \qquad \forall v_h \in V_h.
\end{equation*}
Equivalently, in terms of coefficients
\begin{equation*}
    ( \mat M + \tau \theta \mat K ) \vec u^{n+1} = (\mat M - \tau (1-\theta) \mat K) \vec u^n + \tau \Big( \theta \vec b((n+1)\tau) + (1-\theta) \vec b(n \tau) \Big).
\end{equation*}
The following stability and consistency argument gives explicit time-error bounds; see also \cite[Chapter 11]{QuarteroniValli1994} for estimates in terms of the data.
\begin{theorem}
    \label{thm:Galerkin parabolic convergence in tau}
    Assume that $a:V\times V\to\mathbb R$ is bilinear, symmetric, continuous and coercive, that $f\in H^1(0,T;H)$, and that $u_h^0=u_h(0)$. Define
    \begin{equation*}
        \lambda_{\max,h}=\max_{v_h\ne0}\frac{a(v_h,v_h)}{\|v_h\|_H^2}
        =\lambda_{\max}(\mat M^{-1}\mat K).
    \end{equation*}
    For $0\le\theta<\frac12$, assume the CFL condition
    \begin{equation*}
        \tau(1-2\theta)\lambda_{\max,h}\le2.
    \end{equation*}
    For $\frac12\le\theta\le1$ no such restriction is needed. Then
    \begin{equation*}
        \max_{n\tau\le T}\|u_h^n-u_h(n\tau)\|_H
        \le \max\{\theta,1-\theta\}\,\tau\int_0^T\|u_h''(s)\|_H\diff s.
    \end{equation*}
    If $\theta=\frac12$ and $f\in H^2(0,T;H)$, then
    \begin{equation*}
        \max_{n\tau\le T}\|u_h^n-u_h(n\tau)\|_H
        \le \frac{\tau^2}{8}\int_0^T\|u_h^{(3)}(s)\|_H\diff s.
    \end{equation*}
    These are first- and second-order estimates for fixed $h$. Uniform estimates as $h\to0$ require uniform bounds on the displayed derivative integrals.
\end{theorem}
\begin{proof}
    Let $A_h:V_h\to V_h$ satisfy $(A_hv_h,w_h)_H=a(v_h,w_h)$. In an $H$-orthonormal eigenbasis its eigenvalues are positive. The amplification factors of the method are
    \begin{equation*}
        R_\theta(\tau\lambda)=\frac{1-(1-\theta)\tau\lambda}{1+\theta\tau\lambda}.
    \end{equation*}
    Under the stated restriction their absolute values are at most one, and $\|(I+\tau\theta A_h)^{-1}\|_{H\to H}\le1$.
    Substituting the exact semidiscrete solution into a time step gives the residual
    \begin{equation*}
        r_n=\frac1\tau\int_{n\tau}^{(n+1)\tau}u_h'(s)\diff s
        -(1-\theta)u_h'(n\tau)-\theta u_h'((n+1)\tau).
    \end{equation*}
    The fundamental theorem of calculus bounds $\|r_n\|_H$ by
    $\max\{\theta,1-\theta\}\int_{n\tau}^{(n+1)\tau}\|u_h''(s)\|_H\diff s$.
    For $\theta=\frac12$, the trapezoidal-rule remainder instead gives
    $\|r_n\|_H\le(\tau/8)\int_{n\tau}^{(n+1)\tau}\|u_h^{(3)}(s)\|_H\diff s$.
    The contraction estimate for each step yields an error bounded by $\tau\sum_n\|r_n\|_H$, proving both claims.
\end{proof}

\begin{remark}
    If $|a(v,w)|\le M_a\|v\|_V\|w\|_V$, then $\lambda_{\max,h}\le M_a C(h)$, where $C(h)=\sup_{v_h\ne0}\|v_h\|_V^2/\|v_h\|_H^2$. Thus $\tau(1-2\theta)M_aC(h)\le2$ is a sufficient CFL restriction when $\theta<\frac12$. A strict safety margin gives damping of the highest-frequency modes.
    Combining the spatial and temporal bounds by the triangle inequality estimates $\|u(n\tau)-u_h^n\|_H$.
\end{remark}

\subsection*{Exercises}
\begin{exercises}
    \ex Implement the explicit Euler method for the heat equation $\frac{\partial u}{\partial t} - \frac{\partial^2 u}{\partial x^2} = f$ in the interval $(0,1)$ with homogeneous Dirichlet boundary conditions using piece-wise linear basis functions. Avoid using any numerical library for ODEs.
    Compute the convergence error graph for the case where $f = 0$ and initial condition $u_0(x) = \sin(\pi x)$, where the exact solution is given by $u(x,t) = e^{-\pi^2 t} \sin(\pi x)$.
\end{exercises}
